% Generated draft: requires mathematical and source review.
\documentclass[11pt]{article}
\usepackage[margin=25mm]{geometry}
\usepackage{fontspec}
\setmainfont{Latin Modern Roman}
\usepackage{amsmath,amssymb,mathtools,unicode-math}
\setmathfont{Latin Modern Math}
\usepackage{enumitem,needspace,xurl,hyperref}
\setlength{\parindent}{0pt}
\setlength{\parskip}{5pt}
\setlength{\emergencystretch}{3em}
\newcommand{\N}{\mathbb{N}}
\newcommand{\Z}{\mathbb{Z}}
\newcommand{\Q}{\mathbb{Q}}
\newcommand{\R}{\mathbb{R}}
\newcommand{\C}{\mathbb{C}}
\newcommand{\F}{\mathbb{F}}
\newcommand{\A}{\mathbb{A}}
\newcommand{\PP}{\mathbb{P}}
\newcommand{\E}{\mathbb{E}}
\newcommand{\eps}{\varepsilon}
\newcommand{\dd}{\,\mathrm{d}}
\begin{document}
\section*{UnsolvedMath ID: 800908}
\textbf{Source problem code:} catalogue entry
\subsection{Definitions and mathematical statement}
Use the standard pluripotential interpretation of the source terminology: \(X\) is a compact K\"ahler manifold, and \(\omega\) is a smooth, real, \(d\)-closed \((1,1)\)-form. Semipositivity means \(\omega\geq 0\) pointwise. Bigness means that there exist a K\"ahler form \(\eta\) on \(X\), a number \(\varepsilon>0\), and a quasi-plurisubharmonic function \(\psi\) such that
\[
\omega+i\partial\bar\partial\psi\geq\varepsilon\eta
\]
in the sense of currents. A quasi-plurisubharmonic function is locally the sum of a plurisubharmonic function and a smooth real function. These standard conventions for the type, closedness, and bigness of \(\omega\) are not explicitly expanded in the supplied source excerpt.

Let \(PSH(X,\omega)\) denote the upper semicontinuous, locally integrable quasi-plurisubharmonic functions \(\varphi:X\to[-\infty,\infty)\) satisfying
\[
\omega+i\partial\bar\partial\varphi\geq0
\]
in the sense of currents. A smooth \(\omega\)-plurisubharmonic function is a function \(u\in C^\infty(X,\R)\) satisfying \(\omega+i\partial\bar\partial u\geq0\) pointwise.

\textbf{Question 39 (attributed in the source to Eyssidieux--Guedj--Zeriahi).} For every compact K\"ahler manifold \(X\), every smooth semipositive and big form \(\omega\) as above, and every \(\varphi\in PSH(X,\omega)\), does there exist a sequence \((\varphi_j)_{j\geq1}\subset C^\infty(X,\R)\) such that
\[
\omega+i\partial\bar\partial\varphi_j\geq0,\qquad
\varphi_{j+1}(x)\leq\varphi_j(x)
\quad(j\geq1,\ x\in X),
\]
and
\[
\lim_{j\to\infty}\varphi_j(x)=\varphi(x)
\quad\text{for every }x\in X?
\]
The limit includes points where \(\varphi(x)=-\infty\). The source does not specify its convention concerning identically \(-\infty\) functions; the locally integrable convention used here excludes them on each connected component.
\subsection{Short English statement}
Does every \(\omega\)-plurisubharmonic function on a compact K\"ahler manifold, with \(\omega\) smooth, semipositive, and big, admit a pointwise decreasing approximation by globally smooth functions that remain \(\omega\)-plurisubharmonic? The supplied registry labels the question open; this status has not been independently verified.
\subsection{Sources}
\begin{itemize}
\item Supplied local UnsolvedMath registry record: permanent ID \(800908\), exact source code \texttt{catalogue entry}, title \emph{Open Problems in Pluripotential Theory --- Question 39}. Supplied upstream revision: \texttt{d371ee917ac9647fbb06cd6fc529deec6f3dda82}.
\item Dinew, Guedj, and Zeriahi, \emph{Open Problems in Pluripotential Theory} (2015), Question 39. Supplied source URL: \url{https://arxiv.org/abs/1511.00705}. The registry lists Dinew, Guedj, and Zeriahi as proposers and gives the question-specific attribution \textnormal{[Eyssidieux--Guedj--Zeriahi]}.
\item Supplied provenance records access on \texttt{2026-10-16} and extraction method \texttt{arxiv\_latex}. These are supplied provenance details, not a claim of independent source access or verification. The supplied status is \texttt{open}.
\end{itemize}
\end{document}
